%=================================
\chapter{MAIN 17: MẪU SOẠN CHUYÊN ĐỀ}
%==================BÀI SỐ 1
%\setcounter{section}{123}
\section{Bất phương trình bậc hai một ẩn}
\subsection{Tóm tắt lí thuyết}
\begin{tomtat}
\subsubsection{Phần nội dung thứ nhất}
	\begin{dn}
		Tam thức bậc hai là biểu thức có dạng $f(x)  =  ax^2  +  bx  +  c,$ trong đó $a,b,c$ là những hệ số, $a\ne 0$. Nghiệm của tam thức bậc hai là giá trị của $x$ làm cho tam thức có giá trị bằng $0$.
	\end{dn}	
	\begin{dl}
		Toanvip307.blogspot.com ---
		Cung cấp code \LaTeX{} cái loại bìa sách, logo, mẫu sách, .... giá rẻ. Nhận thiết kế bản quyền và cover mọi loại bìa sách, logo, mẫu sách (trong khả năng).
		\begin{listEX}[2]
			\item FB: Trần Thịnh.
			\item Số ĐT: 09.7828.7121
		\end{listEX}
	\end{dl}
\lipsum[1] %lấy văn bản test Main
\begin{note}
	Cung cấp code \LaTeX{} cái loại bìa sách, logo, mẫu sách, .... giá rẻ. Nhận thiết kế bản quyền và cover mọi loại bìa sách, logo, mẫu sách (trong khả năng).
\end{note}
\begin{hq}
	Cung cấp code \LaTeX{} cái loại bìa sách, logo, mẫu sách, .... giá rẻ. Nhận thiết kế bản quyền và cover mọi loại bìa sách, logo, mẫu sách (trong khả năng).
\end{hq}
\begin{noidung}{MẪU KHUNG NÀY VẪN CÒN}
	Làm bìa sách, logo, main, .... vất xà vả lắm thầy cô. Giá chỉ vài lon bia, lại được cập nhật hỗ trợ nhiệt xà tình.
	\begin{enumerate}[\itemKN]
	\item FB: Trần Thịnh.
	\item Zalo: 09.7828.7121
	\item Đảm bản uy tín chất lượng.
	\end{enumerate}
\end{noidung}
\subsubsection{Giới thiệu môt số môi trường khác dùng trong main}
\lipsum[1]
\begin{note}
	{\bfseries Toanvip307.blogspot.com}\\
	Toanvip307.blogspot.com cung cấp code \LaTeX{} cái loại bìa sách, logo, mẫu sách, .... giá rẻ. Nhận thiết kế bản quyền và cover mọi loại bìa sách, logo, mẫu sách
	\begin{listEX}[2]
		\item Fb: Trần Thịnh
		\item ĐT: 09.7828.7121
	\end{listEX}
\end{note}

\begin{dl}
	Cho tam thức bậc hai $f(x)  =  ax^2  +  bx  +  c,a\ne 0,\Delta  =  b^2  -  4ac$.
	Khi đó:
	\begin{itemize}
		\item $\Delta  <  0\Rightarrow af(x)  >  0,\forall x\in \mathbb{R}$.
		\item $ \Delta  =  0\Rightarrow af(x)  >  0,\forall x\in \mathbb{R}\backslash \left\{ -  \dfrac{b}{2a}\right\} $ và $f\left(-\dfrac{b}{2a}\right)=0$.
		\item $ \Delta  >  0\Rightarrow \heva{
			& af(x)  >  0,\,\forall x\in (-\infty; x_1)\cup (x_2; +\infty) \\ 
			& af(x)  <  0,\,\forall x\in (x_1; x_2)}$ (Với $x_1<x_2$ là nghiệm của $f(x)$).
		\begin{center}
			\begin{tikzpicture}
				\tkzTabInit[nocadre=false, lgt=1, espcl=2.75]{$x$ /0.75,$f(x)$ /0.6}{$-\infty$,$x_1$,$x_2$,$+\infty$}
				\tkzTabLine{, $af(x)>0$,0,$af(x)<0$,0,$af(x)>0$,}
			\end{tikzpicture}
		\end{center}
	\end{itemize}	
\end{dl}

\begin{binhluan}\\
	Với $b'=\dfrac{b}{2}$ và $\Delta'  =  b'^2  -  ac$.
	Ta có kết quả tương tự:
	\begin{itemize}
		\item $\Delta'  <  0\Rightarrow af(x)  >  0,\,\forall x\in \mathbb{R}$.
		\item $ \Delta'  =  0\Rightarrow af(x)  >  0,\,\forall\, x\in \mathbb{R}\backslash \left\{ -  \dfrac{b'}{a}\right\} $ và $f\left(-\dfrac{b'}{a}\right)=0$.
		\item $ \Delta'  >  0\Rightarrow \heva{
			& af(x)  >  0,\forall x\in (-\infty; x_1)\cup (x_2; +\infty) \\ 
			& af(x)  <  0,\forall x\in (x_1; x_2)}$
	\end{itemize}	 
\end{binhluan}
\begin{nx}
	Toanvip307.blogspot.com cung cấp code \LaTeX{} cái loại bìa sách, logo, mẫu sách, .... giá rẻ. Nhận thiết kế bản quyền và cover mọi loại bìa sách, logo, mẫu sách
\end{nx}
\begin{luuy}
	Tiếc gì 1 thùng bia Sài Gòn mà không đăng kí sử dụng Main 15 này nhỉ?
	\begin{listEX}[1]
		\item Fb: Trần Thịnh
		\item ĐT: 09.7828.7121
	\end{listEX}
\end{luuy}
\begin{tc}
	Toanvip307.blogspot.com cung cấp code \LaTeX{} cái loại bìa sách, logo, mẫu sách, .... giá rẻ. Nhận thiết kế bản quyền và cover mọi loại bìa sách, logo, mẫu sách
	\begin{listEX}[2]
		\item Fb: Trần Thịnh
		\item ĐT: 09.7828.7121
	\end{listEX}
\end{tc}




\subsubsection{Bất phương trình bậc hai một ẩn}

\begin{dn}
	Bất phương trình bậc hai ẩn $x$ là bất phương trình dạng $ax^2+bx+c>0$ (hoặc $ax^2+bx+c<0$, $ax^2+bx+c\le 0$,  $ax^2+bx+c\ge 0$), trong đó $a,b,c$ là những số thực đã cho, $a\ne 0$.
\end{dn}
\lipsum[4]
\begin{ghichu}
	Thầy Thịnh Toán - Quảng Trị. Dạy Toán lớp 6-12. Ôn thi vào lớp 10. Luyện thi THPT Quốc gia.
\end{ghichu}
\begin{khung4}{Lớp Toán thầy Thịnh}
	Thầy Thịnh Toán - Quảng Trị. Dạy Toán lớp 6-12. Ôn thi vào lớp 10. Luyện thi THPT Quốc gia.
	\begin{listEX}[2]
		\item FB: Trần Thịnh.
		\item Số ĐT: 09.7828.7121
	\end{listEX}
\end{khung4}
\lipsum[4]
\begin{noidung}{GHI NHỚ 
	}
	Thầy Thịnh Toán - Quảng Trị. Dạy Toán lớp 6-12. Ôn thi vào lớp 10. Luyện thi THPT Quốc gia.
	\begin{enumerate}[\itemCI]
		\item FB: Trần Thịnh.
		\item Số ĐT: 09.7828.7121
		\item! Địa chỉ: Triệu Thành, Triệu Phong, Quảng Trị
	\end{enumerate}
\end{noidung}
\end{tomtat}


\subsection{Các dạng toán thường gặp}
\begin{dang}{Xét tính đơn điệu của hàm số cho bởi công thức}
	\begin{itemize}
		\item \textit{Bước 1.} Tìm tập xác định $\mathscr{D}$ của hàm số.
		\item \textit{Bước 2.} Tính đạo hàm $f'(x)$ của hàm số. Tìm các điểm $x_1$; $x_2$;\ldots; $x_n$ thuộc $\mathscr{D}$ mà tại đó đạo hàm $f'(x)$ bằng $0$ hoặc không tồn tại.
		\item \textit{Bước 3.} Sắp xếp các điểm $x_1$; $x_2$;\ldots; $x_n$ theo thứ tự tăng dần, xét dấu $f'(x)$ và lập bảng biến thiên.
		\item \textit{Bước 4.} Nêu kết luận về các khoảng đồng biến, nghịch biến của hàm số.
	\end{itemize}
\end{dang}
\begin{vd}%[BG12new-4in1, Trần Hoà]%[[2D1H1-1]
	Tim khoảng đồng biến, nghịch biến của hàm số $y=-2x^2+4x+3$.
	\loigiai{
		Hàm số đã cho có tập xác định là $\mathbb{R}$.\\
		Ta có: $y'=-4x+4$;\\
		$y'=0\Leftrightarrow -4x+4=0\Leftrightarrow x=1$.\\
		Ta có bảng xét dấu $y'$ như sau:
		\begin{center}
			\begin{tikzpicture}[font=\footnotesize,thick,>=stealth]
				\tikzset{double style/.append style={double distance=1.5pt}}
				\tkzTabInit[nocadre=false,lgt=1.2,espcl=2.5,deltacl=0.6]
				{$x$ /0.6, $y'$ /0.6}
				{$-\infty$,$1$,$+\infty$}
				\tkzTabLine{ ,+,$0$,-, }
			\end{tikzpicture}
		\end{center}
		Vậy hàm số đồng biến trên khoảng $(-\infty;1)$; nghịch biến trên khoảng $(1;+\infty)$.
	}
\end{vd}
\begin{vd}%[BG12new-4in1, Trần Hoà]%[2D1V1-1]
	Cho hàm số $y=x\mathrm{e}^x$. Mỗi khẳng định sau đúng hay sai?
	\choiceTF[t]
	{$y'=\mathrm{e}^x$}
	{\True Hàm số đồng biến trên khoảng $(-1;+\infty)$}
	{$y+y'=y''$}
	{\True $y\left(2^{2023}\right)<y\left(2^{2024} \right)$}
	\loigiai{
		
		\begin{itemchoice}
			\itemch Sai. Ta có $y'=x'\cdot \mathrm{e}^x+(\mathrm{e}^x)'\cdot x=\mathrm{e}^x+x\mathrm{e}^x=(1+x)\mathrm{e}^x$.
			\itemch Đúng. $y'=(1+x)\mathrm{e}^x$. Có $y'>0\Leftrightarrow (1+x)\mathrm{e}^x>0\Leftrightarrow 1+x>0\Leftrightarrow x>-1$.\\
			Vậy hàm số đồng biến trên khoảng $(-1;+\infty)$.
			\itemch Sai. Ta có $y''=(1+x)'\cdot\mathrm{e}^x+\left(\mathrm{e}^x\right)'\cdot (1+x)=\mathrm{e}^x +(1+x)\mathrm{e}^x=(2+x)\mathrm{e}^x$.\\
			Lại có $y+y'=x\mathrm{e}^x+(1+x)\mathrm{e}^x=(1+2x)\mathrm{e}^x$. Suy ra $y+y'\neq y''$.
			\itemch Đúng. Vì hàm số đồng biến trên khoảng $(-1;+\infty)$ và $-1<2^{2023}<2^{2024}$ nên $y\left(2^{2023}\right)<y\left(2^{2024} \right)$.
		\end{itemchoice}
	}	
\end{vd}
\begin{vd}%[BG12new-4in1, Trần Hoà]%[2D1V1-1]
	Cho hàm số $y=f(x)$ có đạo hàm $f'(x)=(1-x)x$. Hỏi khoảng nghịch biến của hàm số $y=f(3-2x)$ có độ dài lớn nhắt bằng bao nhiêu?
	\shortans[]{$0{,}5$}
	\loigiai{
		Ta có $f'(x)=0\Rightarrow\hoac{&x=1\\&x=0.}$\\
		Bảng xét dấu
		\begin{center}
			\begin{tikzpicture}\tkzTabInit[nocadre=false,lgt=1.2,espcl=2.5,deltacl=0.6]
				{$x$ /0.6,$f'$ /0.6}
				{$-\infty$,$0$,$1$,$+\infty$}
				\tkzTabLine{,-,$0$,+,$0$,-,}
			\end{tikzpicture}
		\end{center}
		Ta có $y=f(3-2x)\Rightarrow y'=-2f'(3-2x)$.\\
		Từ bảng bảng xét dấu $y'=f'(x)$, ta có hàm số $y=f(3-2x)$ nghịch biến khi\\ $y'=-2f'(3-2x)<0\Leftrightarrow f'(3-2x)>0\Leftrightarrow 0<3-2x<1\Leftrightarrow \dfrac{3}{2}>x>1$.\\
		Vậy hàm số nghịch biến trên khoảng $\left(1;\dfrac{3}{2}\right)$.\\ Do đó khoảng nghịch biến có độ dài lớn nhất là $\dfrac{3}{2}-1=0{,}5$.
	}
\end{vd}


\subsection{Bài tập}
\subsubsection{Bài tập tự luận}
\begin{bt}%[BG12new-4in1, Trần Hoà]%[2D1H1-1]
	Tìm các khoảng đơn điệu của hàm số $y=-x^3+6x^2-9x+4$.
	\loigiai{
		Hàm số $y=-x^3+6x^2-9x+4$ có tập xác định $\mathscr{D}=\mathbb{R}$.\\
		Ta có $y'=-3x^2+12x-9$. Cho 
		{\allowdisplaybreaks
			\begin{eqnarray*}
				y'&=&0 \Leftrightarrow -3x^2+12x-9=0\\
				&\Leftrightarrow&
				\hoac{&x=1\\&x=3.}
		\end{eqnarray*}}	
		Bảng biến thiên
		\begin{center}
			\begin{tikzpicture}
				\tkzTabInit[lgt=1,espcl=3]
				{$x$/0.7,$y'$/0.7,$y$/2}
				{$-\infty$,$1$,$3$,$+\infty$}
				\tkzTabLine{,-,0,+,0,-,}
				\tkzTabVar{+/$+\infty$,-/ $0$/,+/$4$,-/$-\infty$}
			\end{tikzpicture}
		\end{center}
		Vậy hàm số nghịch biến trên mỗi khoảng $(-\infty;1)$, $(3;+\infty)$ và đồng biến trên khoảng $(1;3)$.	
	}
\end{bt}	
\begin{bt}%[BG12new-4in1, Trần Hoà]%[2D1H1-1]
	Tìm các khoảng đơn điệu của hàm số $y=x^4-6x^2+8x+1$.
	\loigiai{
		Hàm số $y=x^4-6x^2+8x+1$ có tập xác định $\mathscr{D}=\mathbb{R}$.\\
		Ta có $y'=4x^3-12x+8=0=4(x-1)^2(x+2)$.
		{\allowdisplaybreaks
			\begin{eqnarray*}
				y'=0 &\Leftrightarrow& 4(x-1)^2(x+2)=0\\ 
				&\Leftrightarrow&
				\hoac{&x=-2\\&x=1.}
		\end{eqnarray*}}			
		Bảng biến thiên
		\begin{center}
			\begin{tikzpicture}
				\tkzTabInit[lgt=1,espcl=3]
				{$x$/0.7,$y'$/0.7,$y$/2}
				{$-\infty$,$-2$,$1$,$+\infty$}
				\tkzTabLine{,-,0,+,0,+,}
				\tkzTabVar{+/$+\infty$,-/ $-23$/,R,+/$+\infty$}
			\end{tikzpicture}
		\end{center}
		Vậy hàm số nghịch biến trên khoảng $(-\infty;-2)$ và đồng biến trên khoảng $(-2;+\infty)$.
	}
\end{bt}





\subsubsection{Bài tập chọn đáp án đúng}
\Opensolutionfile{ans}[ans/ansTN-Dung]
\begin{ex}%[BG12new-4in1, Trần Hoà]%[2D1H1-1]
	Các khoảng đồng biến của hàm số $y=x^4-8x^2-4$ là
	\choice
	{$(-\infty;-2)$ và $(0;2)$}
	{$(-2;0)$ và $(0;+\infty)$}
	{\True $(-2;0)$ và $(2;+\infty)$}
	{$(-\infty;-2)$ và $(2;+\infty)$}
	\loigiai{Ta có tập xác định của hàm số là $\mathscr{D}=\mathbb{R}$ và $y'=4x^3-16x$ nên $y'=0\Leftrightarrow \hoac{&x=0\\&x=-2\\&x=2.}$\\
		Khi đó $y'\ge 0\Leftrightarrow 4x^3-16x\ge 0\Leftrightarrow \hoac{&x\ge 0\\&-2\le x\le 0.}$\\
		Vậy hàm số đồng biến trên các khoảng $(-2;0)$ và $(2;+\infty)$.
	}
\end{ex}
\begin{ex}%[BG12new-4in1, Trần Hoà]%[2D1H1-1]
	Cho hàm số $y=f(x)$ có đạo hàm $f'(x)=x^2(x^2-1),\,\forall x\in \mathbb{R}$. Hàm số $y=2f(-x)$ đồng biến trên khoảng
	\choice
	{$\left(2;+\infty\right)$}
	{$\left(-\infty;-1\right)$}
	{$(0;2)$}
	{\True $\left(-1;1\right)$}
	\loigiai{
		Hàm số $y=g(x)=2f(-x)$ có đạo hàm $y'=-2f'(-x)=-2x^2(x^2-1)$. Suy ra
		$$y'=0\Leftrightarrow -2x^2(x^2-1)=0\Leftrightarrow \hoac{&x=0\\&x=\pm 1.}$$
		Bảng biến thiên của hàm số $y=g(x)=2f(-x)$ là
		\begin{center}
			\begin{tikzpicture}[scale=1, font=\footnotesize, line join=round, line cap=round, >=stealth]
				\tkzTabInit[nocadre=false,lgt=1.5,espcl=1.5,deltacl=0.6]
				{$x$/.7,$g'(x)$/.7,$g(x)$/1.5}
				{$-\infty$,$-1$,$0$,$1$,$+\infty$}
				\tkzTabLine{,-,z,+,z,+,z,-,} 
				\tkzTabVar{+/,-/,R,+/,-/}
			\end{tikzpicture}
		\end{center}
		Vậy suy ra hàm số $y=g(x)=2f(-x)$ đồng biến trên khoảng $\left(-1;1\right)$.
	}
\end{ex}
\Closesolutionfile{ans}





\subsubsection{Bài tập đúng sai}
\Opensolutionfile{ans}[ans/ansTN-DungSai]
\begin{ex}%[BG12new-4in1, Trần Hoà]%[2D1V1-1]
	Cho hàm số $y=\ln (1+x)$. Xét tính đúng sai của mỗi mệnh đề sau
	\choiceTF[t]
	{Hàm số có tập xác định là $\mathbb{R}$}
	{\True Hàm số đồng biến trên tập xác định}
	{\True Tiếp tuyến của đồ thị vuông góc với đường phân giác góc phần tư thứ hai cắt hai trục toạ độ tại hai điểm phân biệt}
	{$\dfrac{1}{y'(0)}+\dfrac{1}{y'(1)}+\dfrac{1}{y'(2)}+\ldots+\dfrac{1}{y'(2024)}=2025\cdot 2026$}
	\loigiai{
		
		\begin{itemchoice}
			\itemch Sai. Tập xác định $\mathscr{D}=(-1;+\infty)$.
			\itemch Đúng. Ta có $y'=\dfrac{(1+x)'}{(1+x)}=\dfrac{1}{1+x}>0,\,\forall x\in (-1;+\infty)$ nên hàm số đồng biến trên tập xác định.
			\itemch Đúng. Đường phân giác góc phần tư thứ hai có phương trình $y=-x$.\\
			Gọi $k$ là hệ số góc của tiếp tuyến cần tìm, ta có $k=1$.\\
			Hoành độ tiếp điểm là nghiệm phương trình $y'=0\Leftrightarrow \dfrac{1}{1+x}=1\Rightarrow x=0\Rightarrow A(0;1)$.\\
			Phương trình tiếp tuyến là $y=1(x-0)+1\Leftrightarrow y=x+1$.\\
			Đường thẳng cắt hai trục toạ độ tại hai điểm phân biệt.
			\itemch Sai. Ta có
			{\allowdisplaybreaks
				\begin{eqnarray*}
					&&\dfrac{1}{y'(0)}+\dfrac{1}{y'(1)}+\dfrac{1}{y'(2)}+\ldots+\dfrac{1}{y'(2024)}\\
					&=&
					1+2+3+\ldots +2025=\dfrac{2025\cdot 2026}{2}.
			\end{eqnarray*}}	
			
		\end{itemchoice}
	}	
\end{ex}
\begin{ex}%[BG12new-4in1, Trần Hoà]%[2D1V1-1]
	Cho hàm số $y=x+\sqrt{4-x^2}$. Xét tính đúng sai của mỗi mệnh đề sau
	\choiceTF[t]
	{\True Hàm số có tập xác định là $\mathscr{D}=[-2;2]$}
	{$y'=1+\dfrac{x}{\sqrt{4-x^2}}$}
	{Khoảng đồng biến trên khoảng $(-2;2)$}
	{Phương trình $x+\sqrt{4-x^2}=3$ có hai nghiệm phân biệt}
	\loigiai{
		
		\begin{itemchoice}
			\itemch Đúng. Điều kiện: $4-x^2\ge 0\Leftrightarrow -2\le x\le 2 $. Tập xác định $\mathscr{D}=[-2;2]$.
			\itemch Sai. Ta có $y'=1+\dfrac{(4-x^2)'}{2\sqrt{4-x^2}}=1-\dfrac{x}{\sqrt{4-x^2}}$.
			\itemch Sai. Ta có $y'=0\Leftrightarrow \dfrac{x}{\sqrt{4-x^2}}=1\Leftrightarrow x=\sqrt{4-x^2}\Leftrightarrow\heva{&x\ge 0\\&x=\pm \sqrt{2}}\Leftrightarrow x=\sqrt{2}.$\\
			Bảng xét dấu
			\begin{center}
				\begin{tikzpicture}
					\tkzTabInit[nocadre=false,lgt=1,espcl=2.5,deltacl=0.6]%[nocadre=true\fasle(không\có) kẻ viền bảng ,
					%lgt=1,espcl=2]:Độ rộng cột thứ nhât, các cột tiếp theo
					{$x$ /.7,$y'$ /.7,$y$ /2}{$-2$,$\sqrt{2}$,$2$}
					%số dòng của x(1), y'(1), y(2), các giá trị điền ở dòng
					\tkzTabLine{,+,$0$,-,}%điền dấu, giá trị y'
					\tkzTabVar{-/ $0$ ,+/$2\sqrt{2}$,-/$0$} %-/ là giá trị ở dưới thấp
				\end{tikzpicture}
			\end{center}
			Hàm số đồng biến trên khoảng $(-2;\sqrt{2})$ và nghịch biến trên khoảng $(\sqrt{2};2)$. 
			\itemch Sai. Từ bảng biến thiên ta thấy đường thẳng $y=3$ không cắt đồ thị nên phương trình vô nghiệm. 		
		\end{itemchoice}
	}	
\end{ex}
\Closesolutionfile{ans}






\subsubsection{Bài tập điền khuyết}
\Opensolutionfile{ans}[ans/ansTN-DienKhuyet]
\begin{ex}%[BG12new-4in1, Trần Hoà]%[2D1H1-1]
	Hàm số $y=x^4+4x^2+3$. Có bao nhiêu giá trị $x$ nguyên thuộc khoảng đồng biến của hàm số và $x\in [-10;10]$?
	\shortans[]{$-9$}
	\loigiai{
		Ta có $y'=4x^3+8x=4x(x^2+2)$; $y'=0 \Leftrightarrow x=0$.\\
		Bảng biến thiên:
		\begin{center}
			\begin{tikzpicture}
				\tkzTabInit[nocadre=false,lgt=1,espcl=2.5]
				{$x$ /0.6,$y'$ /0.6,$y$ /1.7}
				{$-\infty$,$0$,$+\infty$}
				\tkzTabLine{,-,$0$,+,}
				\tkzTabVar{+/$+\infty$, -/$3$,+/$+\infty$}
			\end{tikzpicture}
		\end{center}
		Suy ra hàm số đồng biến trên khoảng $(0;+\infty )$.\\
		Kết hợp với điều kiện $x$ nguyên thuộc đoạn $[-10;10]$, ta có $x\in\{1;2;\ldots;10\}$.
	}
\end{ex}
\begin{ex}%[BG12new-4in1, Trần Hoà]%[2D1H1-1]
	Hàm số $y=x^4+4x^2+3$. Có bao nhiêu giá trị $x$ nguyên thuộc khoảng đồng biến của hàm số và $x\in [-10;10]$?
	\shortans[]{$-9$}
	\loigiai{
		Ta có $y'=4x^3+8x=4x(x^2+2)$; $y'=0 \Leftrightarrow x=0$.\\
		Bảng biến thiên:
		\begin{center}
			\begin{tikzpicture}
				\tkzTabInit[nocadre=false,lgt=1,espcl=2.5]
				{$x$ /0.6,$y'$ /0.6,$y$ /1.7}
				{$-\infty$,$0$,$+\infty$}
				\tkzTabLine{,-,$0$,+,}
				\tkzTabVar{+/$+\infty$, -/$3$,+/$+\infty$}
			\end{tikzpicture}
		\end{center}
		Suy ra hàm số đồng biến trên khoảng $(0;+\infty )$.\\
		Kết hợp với điều kiện $x$ nguyên thuộc đoạn $[-10;10]$, ta có $x\in\{1;2;\ldots;10\}$.
	}
\end{ex}
\Closesolutionfile{ans}


\newpage
\indapanT{10}{ans/ansTN-Dung}{BÀI TẬP CHỌN ĐÁP ÁN}
\indapanTF{3}{ans/ansTN-DungSai}{BÀI TẬP ĐÚNG SAI}
\indapanSA{6}{ans/ansTN-DienKhuyet}{BÀI TẬP ĐIỀN KHUYẾT}